\section{Approximations as Tangents}
\label{sec:approx-as-tangents}

We motivate our approach by showing how to combine ideas from differential geometry and linear algebra to
reconstruct the dependency analyses of \secref{introduction} in a denotational setting.

\subsection{Manifolds, Smooth Functions, and Automatic Differentiation}
\label{sec:approx-as-tangents:autodiff}

% For what follows we do not need to know in detail what a manifold is, or exactly how smooth functions are defined. We only state some definitions in order to fix terminology and to justify our claim of a connection to manifolds and smooth functions. \bob{Need a good reference here.}

The general study of differentiable functions takes place on \emph{manifolds}, topological spaces that
``locally'' behave like an open subset of the Euclidean space $\RR^n$. The spaces $\RR^n$ themselves are
manifolds, but so are ``non-flat'' examples such as $n$-spheres and yet more exotic spaces. Every point $x$ in
a manifold $M$ has an associated \emph{tangent vector space} $\tangents_x(M)$ consisting of linear
approximations of curves on the manifold passing through $x$. Each point also has a \emph{cotangent vector
space} $\cotangents_x(M) = \tangents_x(M) \linearto \RR$. The tangent and cotangent spaces are finite
dimensional, so in the presence of a chosen basis they are canonically isomorphic. In the case when the
manifold is $\RR^n$, then every tangent space is isomorphic to $\RR^n$ as well.

Smooth functions $f$ between manifolds $M$ and $N$ are functions on their points that are locally
differentiable on $\RR^n$. Manifolds and smooth functions form a category $\Man$. Each smooth function induces
maps of the (co)tangent spaces:
\begin{itemize}[leftmargin=\enummargin]
\item The \emph{forward derivative} (tangent map, pushforward) $\pushf{f}_x$ is a linear map $\tangents_x(M) \linearto \tangents_{f(x)}(N)$. In the Euclidean case when $M = \RR^m$ and $N = \RR^n$, the tangent map can be represented by the Jacobian matrix of partial derivatives of $f$ at $x$.
\item The \emph{backward derivative} (cotangent map, pullback) $\pullf{f}_x$ is a linear map $\cotangents_{f(x)}(N) \linearto \cotangents_x(M)$. In the Euclidean case, the backward derivative is represented by the transpose of the Jacobian of $f$ at $x$.
\end{itemize}

\begin{remark}[Chain Rule]
  \label{rem:chain-rule}
  A useful property of derivative maps is that they compose according
  to the chain rule. Suppose that $f : M \to N$ and $g : N \to K$ are
  smooth functions. Then for any $x \in M$, we have:
  \begin{itemize}
  \item $\pushf{(g \circ f)}_x = \pushf{g}_{f(x)} \circ \pushf{f}_x : \tangents_x(M) \linearto \tangents_{g(f(x))}(K)$
  \item $\pullf{(g \circ f)}_x = \pullf{f}_x \circ \pullf{g}_{f(x)} : \cotangents_{g(f(x))}(K) \linearto \cotangents_x(M)$
  \end{itemize}
  The chain rule has the practical effect that we can compute
  derivative maps of $f$ and $g$ independently and compose them,
  instead of the potentially more difficult task of computing the
  derivative maps of $g \circ f$. As we shall see below, stable maps
  also obey a chain rule, and this forms the basis of the general
  categorical approach to differentiability that we describe in
  \secref{models-of-total-gps}.
\end{remark}

Computing the forward and backward derivatives of smooth functions $f$ has many applications of practical
interest. For example, computation of the reverse derivative is of central interest in machine learning by
gradient descent, the main technique used to train deep neural networks~\cite{rumelhart88,goodfellow16}.

Derivatives can be computed numerically by computing $f$ on small perturbations of its input, or symbolically
by examining a closed-form representation of $f$. However, a more common and practical technique is to use
\emph{automatic differentiation}, where a program computing $f$ is instrumented to produce (a representation
of) the forward and/or backward derivative as a side-effect of producing the output~\cite{linnainmaa76}. This
has led to the area of differentiable programming, where programming languages and their implementations are
specifically designed to admit efficient automatic differentation
algorithms~\cite{jax2018github,abadi16,elliott17,sigal24}.

\subsection{\added{Potential Change as Tangents}}

\Added{Differential Calculus tracks how functions adapt to changes in
  their inputs. The tangent space at a point is the vector space of
  all potential changes applicable at that point. The forward and
  reverse derivatives transport changes forward and back. We track
  data provenance using a simplification of this idea where tangents
  are taken from the Boolean semiring $\{0,1\}$ where $0$ represents
  ``no change'' and $1$ represents ``possible change''. Potential
  change to an $n$-tuple of data is represented as an $n$-vector of
  $\{0,1\}$ values. These vectors form a semimodule over the semiring
  structure on $\{0,1\}$ with maximum (join) as addition and minimum
  (meet) as multiplication. Thus ``no change'' is represented as the
  $0$ vector $(0 \dots 0)$.}

\Added{With this definition of ``tangent'', the derivative of a
  function
  $f : A_1 \times \cdots \times A_m \to B_1 \times \cdots \times B_n$
  ought to map tangent vectors in $\{0,1\}^m$ to tangent vectors in
  $\{0,1\}^n$, describing how potential changes in the input are
  reflected in the output. Differentiable functions on manifolds have
  an intrinsic notion of derivative, but for plain functions $f$ we
  require that they are paired with a tangent map that describes the
  intensional information about how they depend on their
  arguments. For $f$, the tangent map $\partial f$ has type
  $A_1 \times \cdots \times A_n \to \{0,1\}^{m\times n}$, mapping
  points to matrices representing linear maps: the ``Jacobian'' for
  this function. We are not free to choose an arbitrary linear map for
  $\partial f(a_1, \dots, a_m)$, but must choose one that satisfies
  the following \emph{dependency correctness} property:}

\begin{AddedBlock}
  \begin{definition}[Dependency Correctness] A tangent map
    $\partial f : \overrightarrow{A_i} \to \{0,1\}^{m \times n}$ is
    \emph{dependency correct} for a function
    $f : \overrightarrow{A_i} \to \overrightarrow{B_j}$ if for all
    $a, a' \in \overrightarrow{A_i}$ and $0 \leq j < m$, if $a$ and
    $a'$ are equal on all components $i$ such that
    $\partial f(a)_{i,j} = 1$, then $f(a)_j = f(a')_j$.
  \end{definition}
\end{AddedBlock}

\Added{Reading the definition backwards, this says that $f$ is
  constant at component $j$ on inputs that differ only at components
  $i$ where the tangent map matrix says $0$. Thus we read the entry at
  $(i,j)$ of the tangent matrix as capturing the dependency of output
  $j$ on input $i$. We illustrate this with an example:}

\begin{AddedBlock}
  \begin{example}
    Let $\pi_1 : A_1 \times A_2 \to A_1$ be first projection. At all
    points $(a_1, a_2)$, the tangent matrix is $(1\;0)$ indicating
    that the output depends on the first argument but not the
    second. This is justified by observing how it acts on tangent
    vectors. Representing no change in the first argument and
    potential change in the second as
    $\left(\begin{array}{c} 0 \\ 1 \end{array}\right)$, applying the
    tangent map gives $0$ indicating no change in the output. Swapping
    the potential change to
    $\left(\begin{array}{c} 1 \\ 0 \end{array}\right)$ yields the
    output tangent $1$ indicating that changing the first argument
    changes the output. By linearity, all other potential changes are
    derivable from these. In particular, for all tangent maps, no
    change in the input ($\overrightarrow{0}$) is always mapped to no
    change in the output.
  \end{example}
\end{AddedBlock}

\Added{Note that dependency correct tangent maps for a function are
  not unique. We can always overapproximate the dependency of a
  function on its input. In the case of a binary function like first
  projection, $\partial\pi_1(a_1,a_2) = (1\;1)$ is also permissble.  }

\begin{AddedBlock}
  \begin{example}
    Multiplication $(*) : \mathbb{N} \times \mathbb{N} \to \mathbb{N}$
    is an example of a function that has tangent maps that depend on
    the input. Dependency correct tangent maps can be defined as:
    \begin{displaymath}
      \begin{array}{lcl}
        \partial(*)(0,0)&=&(1\;1) \\
        \partial(*)(0,n)&=&(1\;0) \\
        \partial(*)(m,0)&=&(0\;1) \\
        \partial(*)(m,n)&=&(1\;1)
      \end{array}
    \end{displaymath}
    These maps indicate that the dependency of multiplication on an
    input depends on whether or not other other input is $0$. If the
    first argument is zero, then changing the second will have no
    effect and vice versa. If both arguments are $0$, then a change in
    both will change the output. However, since tangent maps are
    componentwise we have no way of saying change in both, so we
    overapproximate by saying that a change in either \emph{may}
    change the output. Again, this choice is not unique. We could also
    use $(1\;0)$ or $(0\;1)$ and still be dependency correct,
    indicating a choice to be left- or right-biased.
  \end{example}
\end{AddedBlock}

\Added{Functions with dependency correct tangent maps have an identity
  and compose according to the chain rule, so they form a category.}

\newcommand{\DepCat}{\mathbf{DepCat}}

\begin{AddedBlock}
  \begin{proposition}
    The identity tangent map is dependency correct for the identity
    function. Composable functions $f, g$ with dependency correct
    tangent maps $\partial f$ and $\partial g$ are closed under
    composition, with the tangent map being
    $\partial(f \circ g)(x) = \partial f (g(x)) \circ \partial g(x)$.

    The category $\DepCat$ of functions with dependency information
    has objects that are tuples of sets $\overrightarrow{A_i}$ and
    morphisms that are functions with dependency correct tangent
    maps. This category has finite products.
  \end{proposition}

  FIXME: finite products is interesting, because it determines the
  dependency information???

  The fact that dependency correct functions form a category allows us
  to model a simple programming language with finite product
  types. Primitive operations can be interpreted a functions that
  overapproximate their dependency information:

  \begin{proposition}
    There is a functor $\Set \to \DepCat$ that preserves finite
    products. FIXME: check!
  \end{proposition}

  Morphisms in $\DepCat$ constructed from finite product structure and
  embedded functions will only have trivial tangent maps arising from
  the use of projections. A generator of non-trivial dependency
  information is the ability to use parts of the input depending on
  the input itself, via a condition (if-then-else) map:

  \begin{proposition}
    Let $\mathbb{B}$ be the object of $\DepCat$ consisting of the set
    of booleans $\{\mathsf{true}, \mathsf{false}\}$. There are
    dependency correct maps
    $\mathrm{true}, \mathrm{false} : 1 \to \mathbb{B}$ and a natural
    family of maps
    $\mathrm{cond}_X : \mathbb{B} \times (X \times X) \to X$ such that $\mathrm{cond}_X \circ \langle \mathrm{true} \circ !, \mathrm{id} \rangle = \pi_1$ and $\mathrm{cond}_X \circ \langle \mathrm{false} \circ !, \mathrm{id} \rangle = \pi_2$. FIXME: and extensionality.
  \end{proposition}

  Thus $\DepCat$ is suitable for modelling a small language with
  booleans, conditionals, and a collection of primitive types with
  operations. This setting does not cover our example from the
  introduction, which used lists and higher-order functions. In
  Section \ref{sec:models-of-total-gps} we embed $\DepCat$ into
  categories that account for sum types, some algebraic types and
  higher-order functions.
\end{AddedBlock}

\subsection{Derivatives over a Commutative Semiring}
\label{sec:approx-as-tangents:semiring}

In the Euclidean case, the differential picture above is plain matrix algebra: with a chosen basis, tangent
and cotangent spaces are both $\RR^n$, the forward derivative at a point is the Jacobian matrix, and the
backward derivative is its transpose. Once the picture is presented in coordinates like this, the real numbers
play no further role beyond supplying the scalars; since dependency information is typically indexed by
positions in a value, which then serve as a coordinate basis, our first move is to keep the matrix algebra and
change the scalars. Thus throughout this section, we fix a commutative semiring $S = (S, +, 0, \cdot, 1)$,
whose elements we think of as atoms of dependency information, with $\cdot$ combining information along a path
from an input to an output and $+$ combining information across parallel paths.

Our second move is one of abstraction: rather than axiomatise the matrices of \secref{introduction} directly,
we identify the minimal structure required to capture what made that story work: a space of dependency
information for each value, linear maps between spaces for the forward analysis, and an identification of each
space with its dual, the extra structure needed to support an analogue of matrix transposition, and hence a
backward analysis. The role of vector spaces, which are defined over fields, is played by \emph{semimodules}:
commutative monoids $(M, +, 0)$ equipped with an action $a \cdot x$ of the scalars, satisfying the usual laws.
Linear maps preserve $+$, $0$ and the action, and form a category $\SemiMod(S)$; pointwise addition of linear
maps makes each hom-set a commutative monoid, with composition bilinear. Every semimodule has a \emph{dual}
$M^* = M \linearto S$ of linear maps into the scalars (categorically, the internal hom into the scalars
object, which is the unit of the tensor product of semimodules), and every linear map $f : M \linearto N$ has
a \emph{transpose} $\transpose{f} = (- \comp f) : N^* \linearto M^*$, contravariantly.

The transpose is almost the backward analysis we are after: it runs in the opposite direction to $f$, but
between the duals $N^*$ and $M^*$, the cotangent spaces, rather than between the tangent spaces $M$ and $N$
themselves. What we want is a map $N \linearto M$, running dependency information back through the same spaces
the forward map runs it through, as the transposed matrix did in \secref{introduction}. For that, each space
must be identified with its dual. Whereas a finite-dimensional vector space always admits such an
identification (choose a basis), a semimodule need not; we therefore make the identification part of the data.
\begin{definition}[Self-dual semimodule]
  \label{def:self-dual-semimodule}
  A \emph{self-dual} semimodule is a pair $\langle M, i \rangle$ of a semimodule and a chosen isomorphism
  $i : M \cong M^*$. The \emph{pairing} induced by $i$ is $\langle x, y \rangle = i(x)(y)$.
\end{definition}
The pairing is an $S$-valued measure of the extent to which $x$ and $y$ overlap; its prototype is the
dot product of vectors (the pairing induced by the self-duality of $S^n$, as we will see below). Extending the
analogy, we call $x$ and $y$ \emph{orthogonal} when $\langle x, y \rangle = 0$.
\begin{definition}[Conjugate]
  \label{def:conjugate}
  Between self-dual semimodules $\langle M, i \rangle$ and $\langle N, j \rangle$, the \emph{conjugate} of
  a linear map $f : M \linearto N$ is $\conj{f} = i^{-1} \comp \transpose{f} \comp j : N \linearto M$, the
  unique linear map satisfying
  \begin{displaymath}
    \langle f(x), y \rangle = \langle x, \conj{f}(y) \rangle
  \end{displaymath}
\end{definition}
The conjugate is our backward analysis, and its defining equation makes precise the relationship with
the forward one alluded to in \secref{introduction}. Self-dual semimodules and linear maps form a category
$\SDSemiMod(S)$.

The matrices of \secref{introduction} now reappear as concrete instances. The \emph{free} semimodule $S^n$ on
$n$ generators has as elements the $n$-length vectors of scalars, with standard basis vectors $e_i$ and every
$v \in S^n$ decomposing as $v = \sum_i v_i \cdot e_i$. Currying the \emph{dot product} $\langle u, v \rangle =
\sum_i u_i \cdot v_i$ gives an isomorphism $S^n \cong (S^n)^*$, since pairing with the basis vectors extracts
components, $\langle e_i, v \rangle = v_i$; so free semimodules are self-dual. By basis decomposition, a
linear map $S^m \linearto S^n$ is exactly an $n \times m$ matrix over $S$, with composition as matrix
multiplication and the conjugate as the familiar matrix transpose, $(\conj{M})_{ji} = M_{ij}$.

We write $\Mat(S)$ for the skeletal category whose objects are the natural numbers and whose morphisms $m \to
n$ are the $n \times m$ matrices. The Boolean Jacobians from \secref{introduction} live here: a dependency
relation between $m$ input positions and $n$ output positions is exactly a morphism $m \to n$ in $\Mat(\Two)$,
where $\Two$ is the Boolean semiring with $+ = \vee$ and $\cdot = \wedge$, and the chain rule for dependencies
is composition in $\Mat(\Two)$, or equivalently \emph{relational} composition under the equivalence of
categories $\Mat(\Two) \simeq \FinRel$. The dot product can be interpreted as ``overlap'', becauses $\langle
u, v \rangle = \top$ exactly when the selections $u$ and $v$ share a position, i.e.\ are not disjoint.

The trivial semimodule is a zero object in $\SDSemiMod(S)$, and self-dual semimodules are closed under the
\emph{biproduct} $M \biprod N$, an object which is both a product and a coproduct, with the injections
conjugate to the projections, $\biinj_i = \conj{\biproj_i}$. The coincidence of products and coproducts is
characteristic of such linear settings, and means $\SDSemiMod(S)$ on its own cannot interpret programs over
sum types; we return to this in \secref{models-of-total-gps}.

\subsection{Dependency Structures over Semirings}
\label{sec:approx-as-tangents:richer-semirings}

The framework so far only asks for commutativity of the scalars, which makes the dual $M^*$ an $S$-semimodule
again, so that the transpose remains a map between $S$-semimodules. Any commutative semiring $S$ thus yields a
category $\SDSemiMod(S)$ with forward derivatives and their conjugates, related through the pairing. When $S$
is the reals, the conjugate is the familiar adjoint of linear algebra. In a dependency-tracking setting,
however, the relationships of interest are typically granular: we ask whether, or to what degree, an output
depends on an input, and want to compare one answer with another. This calls for an order on dependency
information, which additional axioms on $S$ provide.

\subsubsection{Conjugate pairs and conjugacy analyses}

Suppose first that $1$ is an additive top in $S$ ($1 + a = 1$). Addition is then
idempotent, meaning $a + a = a$. Idempotent addition induces an order: $S$ becomes a join-semilattice,
with $a \leq b$ iff $a + b = b$, joins given by $+$ and least element $0$. This structure lifts pointwise to every semimodule,
and every linear map automatically preserves joins and $0$. Dependency
information is now ordered by informativeness: $0$ means ``no dependency'', and the forward and backward
analyses are monotone, taking more informative inputs to more informative outputs and vice-versa.

If multiplication is also idempotent ($a \cdot a = a$), then $S$ is a bounded distributive lattice,
with multiplication as the meet and distributivity given by the semiring. This opens the way to reading the
two analyses in terms of \emph{disjointness}: where meets exist, define $x \disj y$ when $x \meet y = 0$, and
if disjointness agrees with the orthogonality given by the pairing, the defining equation of the conjugate
(\defref{conjugate}) says that $f$ and $\conj{f}$ exchange disjointness: $f(x) \disj y$ iff $x \disj
\conj{f}(y)$. Let us therefore consider the case where each semimodule has meets, preserved by the scalar action and
agreeing with orthogonality in this sense; meets, unlike joins, do not lift from $S$, so this is additional
structure. Maps that exchange disjointness in this way have a standard name.
\begin{definition}[Conjugate pair]
  \label{def:conjugate-pair}
  Join-preserving maps $f : L \to M$ and $g : M \to L$ between bounded lattices form a \emph{conjugate
  pair} when $f(x) \disj y$ iff $x \disj g(y)$ for all $x \in L$ and $y \in M$.
\end{definition}
The notion originates in relation algebra~\cite{jonsson51}, where the conjugate of composition with a
relation is composition with its converse. The scalars $S$, viewed as a semimodule over itself, qualify
automatically for this reading, and biproducts (including the empty biproduct $S^0$) inherit the structure
componentwise, giving the linear maps between free semimodules $S^n$, along with their conjugates, as the canonical
instances of conjugate pairs.

The relevance to us here is that conjugate pairs also underpin the \emph{cognacy} (common ancestry)
analyses of work on linked visualisations~\cite{perera22,bond25}. These dependency analyses work by composing
a join-preserving dependency analysis with its conjugate, yielding a map sending a selection of outputs to its
\emph{related outputs}, those sharing a dependency on some part of the input; selections in one chart can then
be linked to ``cognate'' selections in another.

\subsubsection{Galois slicing}

Suppose in addition that every scalar has a complement: an operation $\neg$ with $a \meet \neg a =
0$ and $a \join \neg a = 1$, making $S$ a Boolean algebra. Negation supports an alternative presentation of
conjugates as adjoints. Composing a map with complements on either side preserves its direction but exchanges
join-preservation for meet-preservation; the \emph{De Morgan dual} $\neg \comp \conj{f} \comp \neg$ of
the conjugate is thus meet-preserving, and is precisely the upper adjoint of $f$ (upper adjoints
preserve meets, dually to lower adjoints and joins). Forward and backward analysis therefore form a
\emph{Galois connection} between the lattices of approximations.
\begin{definition}[Galois connection]
  \label{def:galois-connection}
  Suppose $X$ and $Y$ are posets. A \emph{Galois connection} $f \adj g: X \to Y$ is a pair of monotone
  functions $f: Y \to X$ and $g: X \to Y$ satisfying $y \leq g(x) \iff f(y) \leq x$ for any $x \in X$ and $y
  \in Y$. Since a Galois connection is also an adjunction, we refer to $f$ as the left adjoint and $g$ as the
  right adjoint.
\end{definition}
This is precisely the setting of \GPS~\cite{perera12a,perera16d,ricciotti17}, recovered here as the
Boolean instance of the general semimodule picture. Galois connections as such presuppose none of the
algebraic structure of this section, being definable between arbitrary posets; De Morgan duality is simply how
the adjoint presentation arises here, where the conjugate is the primary notion. One orientation to keep
straight: the conjugate as presented here, $\conj{f}$, runs backwards, whereas the meet-preserving map of \GPS
runs forwards. The two presentations agree because $\conj{-}$ is involutive: instantiating the construction at
$\conj{f}$ in place of $f$ gives exactly the connection of \GPS, with meet-preserving forward map $\neg \comp
f \comp \neg$.

The join-preserving and meet-preserving maps answer different questions. In the forwards direction, the
join-preserving map $f$ has a \emph{necessity}-flavoured interpretation: which outputs may depend on a given
part of the input. Its meet-preserving counterpart $\neg \comp f \comp \neg$ reads as \emph{sufficiency}:
which outputs a given input selection determines outright. Notably, in the Galois slicing work the
meet-preserving forward map is only used to establish the existence of the join-preserving backward map; it
seems to lack any direct dependency analysis applications of its own.

\subsubsection{Perturbation bounds}
\label{sec:approx-as-tangents:perturbation-bounds}

Dependency analysis over the Booleans records whether an output can change when an input does;
\exref{introduction-example} raised the more fine-grained question of how much it
can change. Numerical analysis studies this question as \emph{propagation of error}: given bounds on how far
each input may be perturbed, derive a bound on the output~\cite{high:ASNA2}. Propagation of error is again a
choice of semiring, this time the \emph{tropical semiring} $(\mathbbm{Q}_{\geq 0} \cup \{\infty\}, \min,
\infty, +, 0)$~\cite{simon88}. A scalar is a non-negative rational (or $\infty$) bounding the extent to which
a value may change from its value in the run, with $\infty$ meaning unconstrained. Semiring addition is
$\min$, so that combining bounds for parallel paths keeps the tighter bound, and multiplication is numerical
$+$, accumulating bounds along a path. Addition is idempotent, so the analyses are monotone as in
\secref{approx-as-tangents:richer-semirings}: smaller bounds are more informative, and the semiring zero
$\infty$ means ``no information''.

A Jacobian entry $c$ now acts as a translation: if the input at that position changes by at most
$\delta$, with the others held fixed, the output changes by at most $\delta + c$. Addition's entries are $0$: a
bound on either argument passes through undisturbed. Multiplication, however, scales perturbations: if $x$
changes by $\delta$ then $x \cdot y$ changes by $|y|\,\delta$, and no translation bounds a scaling; its
entries are $\infty$. This is sound, but reports nothing: propagation of absolute error is blind to
multiplication.

Measuring perturbations multiplicatively instead keeps $\min$ as the join of bounds but accumulates
along paths by numerical product: $(\mathbbm{Q}_{\geq 0}
\cup \{\infty\}, \min, \infty, \times, 1)$, with scalars now read as factors. This is the
multiplicative counterpart of the previous presentation, related to it by logarithms over the reals, and the
two are the classical dichotomy between absolute and relative error~\cite[\S\S 1.6--1.7]{high:ASNA2}.
Multiplication now has unit entries, since relative perturbations pass through a product unchanged, while an
addition $x + y$ acquires the entry $|x| / |x + y|$: the factor by which it amplifies relative error, its
(relative) \emph{condition number}.

In the refund example the final sum adds the products $6$ and $-6$, so the amplification is $|6| / |6 +
(-6)|$: an infinite condition number, known in numerical analysis as catastrophic cancellation. For
non-negative data the factors are at most $1$, so sums never amplify; the refund's negative entry makes the
infinite factor possible. Where the Boolean analysis could only report a dependency that was not really
there, the relative analysis says how badly the bound degrades. The difference between the two is a choice of
semiring, and making that choice is part of specifying what the analysis is to observe.

\subsubsection{Qualitative derivatives}

The Boolean analysis reports a dependency at a cancellation because it forgets signs. The \emph{sign
semiring} $\{+, 0, -, ?\}$, the rule-of-signs domain of abstract interpretation~\cite{cousot77}, retains
them: multiplication of signs is exact, and the only sum whose sign is not determined is $+$ plus $-$, which
is $?$. Taking signs as the scalars gives the \emph{qualitative derivatives} of qualitative
reasoning~\cite{dekleer84,kuipers86}, an entry reading ``increasing this input increases (or decreases, or
cannot affect) that output'', with $?$ appearing exactly at the potential cancellations: in the refund
example, the entry for the price becomes $?$. Collapsing $+$, $-$ and $?$ to $1$ recovers the Boolean
analysis.

\subsubsection{Linearised provenance}
\label{sec:approx-as-tangents:linearised-provenance}

Database provenance has its own semiring tradition: in the framework of
\citet{green07}, query results carry annotations from a commutative semiring, with the \emph{free}
commutative semiring $\mathbbm{N}[X]$ of \emph{provenance polynomials} playing a universal role: its
variables name the input tuples, multiplication records their joint use in a derivation, and addition
collects the alternative derivations. By freeness, assigning a value to each variable extends uniquely to a
semiring homomorphism out of $\mathbbm{N}[X]$, so provenance computed once as polynomials determines, by
evaluation, the annotation in every other semiring.

Provenance semirings can serve as the scalars of the present framework, but with a limitation: a
derivative is linear, so the framework captures \emph{linearised} provenance. Taking $S = \mathbbm{N}[X]$, a
Jacobian entry is a polynomial recording symbolically how an input position contributes to an output position.
Green's semantics assigns an output a polynomial $p$ whose degree records joint use; seeding each input
position with a unique variable, our forward analysis instead returns a polynomial linear in the position
variables, with the partial derivative $\partial p / \partial x_i$, evaluated at the run's input values, as
the coefficient of $x_i$. Joint use survives only in the coefficients: an input
used twice contributes $2x$ where the annotation would record $x^2$.

Translation along a homomorphism applies to the present framework as well: a semiring homomorphism acts
entrywise on Jacobians, and the action is functorial, since matrix composition is built from addition and
multiplication. The translation of interest for dependency tracking is ``zero-testing'' into the Booleans, and
whether that is a homomorphism depends on the source semiring. The counting semiring $\mathbbm{N}$, itself a
provenance semiring recording how many derivations use each input, is \emph{positive}: a sum is zero only when
both summands are, and a product only when a factor is; equivalently, zero-testing is a homomorphism, so
counting translates exactly to dependency tracking. The rationals are not positive, since a number and its
negation sum to zero, and zero-testing $\mathbbm{Q} \to \Two$ preserves multiplication exactly but addition
only laxly. It still acts entrywise on matrices, preserving identities and transposition, but composition only
up to over-approximation: composing Boolean Jacobians can record dependencies absent from the Boolean image of
their composite. The refund in \exref{introduction-example} is a case in point: the price genuinely does not
matter to the total, yet the composed dependencies say it might.
